\section{Background and related work}
\label{sec:background}

\subsection{Leakage in bearing-diagnosis evaluation}

Hendriks, Dumond and Knox showed that the conventional condition-wise split of the CWRU data places the same physical
bearings in training and test sets, and proposed benchmarking with independent sets of bearings \cite{hendriks2022}. Wheat et
al.\ compared run-to-run, day-to-day and part-to-part splits for classical pipelines and found accuracy falling by more than
40\,\% between the most and least separated splits \cite{wheat2024leakage}. Vieira et al.\ generalised bearing-wise
partitioning, reformulated diagnosis as multi-label detection, and measured the inflation caused by segmentation-level and
bearing-level leakage on UORED, Paderborn and CWRU with the training set held fixed and only the test set changed
\cite{vieira2026}. On Paderborn their repetition-wise split --- training and test recordings from the same bearing at the
same condition --- scores near-perfect Macro-AUROC against 54--77\,\% without leakage \cite[Table~12]{vieira2026}. Their study
is supervised and, in their words, concerns ``training and testing within a single testbench''; it does not evaluate
adaptation. Matania et al.\ \cite{matania2024leakage} show how test--training separation should be constructed for bearing
fault classification. Knap et al.\ published a leakage-safe cross-domain benchmark on CWRU and Paderborn with six fixed
scenarios \cite{knap2026leakagesafe}; its default rule separates \emph{recordings}, and one scenario holds out bearings,
which they report among the hardest. Kaya and Jobani report a fall from about 0.86 to 0.56 macro-F1 when Paderborn
bearings are made disjoint in a supervised pipeline \cite{kaya2026}. Concurrently, Nagaswetha and Pathak adapt across
Paderborn conditions with held-out bearings and find that an order-domain representation decides whether alignment helps
\cite{nagaswetha2026}; they use source data and no leaky baseline.

Beyond this field, Kapoor and Narayanan document leakage across 17 scientific disciplines and give the taxonomy we adopt
\cite{kapoor2023leakage}, Saeb et al.\ show the record-wise versus subject-wise analogue in clinical prediction
\cite{saeb2017}, and Geirhos et al.\ name the learning behaviour at stake: a shortcut that succeeds on the benchmark and fails
under the intended shift \cite{geirhos2020shortcut}.

\subsection{Source-free domain adaptation for bearings}

SHOT established the template on which most later methods build: freeze the source classifier, maximise information at the
target and refine centroid-based pseudo-labels \cite{liang2020shot}. SDALR extends it by separating reliable from unreliable
target predictions, using the reliable ones as supervision and the unreliable ones through an entropy term, and reports
\SI{96.8}{\percent} mean accuracy across Paderborn operating conditions \cite{sdalr2025}. Its Paderborn task has eight
classes, each a single physical bearing (K001, KA04, KA15, KA22, KA30, KI14, KI17, KI21), and the same eight specimens form
the source and the target, so the reported transfer is between operating conditions of known bearings. Other source-free
work crosses rigs: Jeong et al.\ combine order normalisation with test-time training and evaluate on a test-only machine
\cite{jeong2025}, and Bio-SFDA reports a CWRU-to-Paderborn task \cite{yoon2026biosfda}. A cross-rig target is
bearing-disjoint by construction, but it changes rig, bearing type, sensor and often label space at once, so it cannot say
how much of the loss is the specimen. Neither design asks how a source-free method fares on unseen bearings of the rig it
was trained on, which is the question here.

\subsection{Physics priors inside adaptation}

Characteristic frequencies computed from geometry and speed are the natural prior for bearing diagnosis. PIUDA builds
physical labels from fault-characteristic spectral energy and aligns them into pseudo-labels for transfer between different machines
\cite{jia2024piuda}. Bio-SFDA gates and weights pseudo-labels with band-energy and signal-to-noise predicates at BPFO, BPFI
and BSF computed on a short-time Fourier transform, with thresholds rescaled online from unlabelled target statistics
\cite{yoon2026biosfda}. PCTL scores the envelope spectrum predicted from the model's features against characteristic
frequencies and uses that score to supervise a confidence predictor while fine-tuning on a small labelled target set
\cite{pctl2026}; it does not filter pseudo-labels. In none of these works is the rate at which the module accepts a fault on
a healthy bearing reported. Because the Bio-SFDA task definition for Paderborn cannot be reconstructed from the article --- a
ball class is listed although the Paderborn real-damage set contains no rolling-element damage --- we evaluate the gating
\emph{mechanism as described, as we implemented it}, and make no comparison against the accuracy reported there.
Schlachter et al.\ simulate pseudo-labelling with ground truth in source-free universal adaptation and report the upper
level that perfect pseudo-labels reach \cite{schlachter2025}; our oracle filter is the keep-or-withhold version of that
instrument.

\subsection{Envelope analysis and the statistics of detection}

The squared envelope spectrum is the standard carrier of second-order cyclostationary bearing signatures \cite{randall2011},
usually after a band selection such as the fast kurtogram \cite{antoni2007kurtogram} and, where discrete components
dominate, after discrete/random separation. Borghesani et al.\ derived the statistics of the squared envelope spectrum under
coloured noise and the corresponding thresholds for testing cyclostationarity \cite{borghesani2013}. The gates compared here
are built from those ingredients; none is claimed as new. Matania et al.\ project order-domain envelope spectra onto
kinematic positions to obtain a machine-invariant representation \cite{matania2025}, the same idea as the harmonic
coordinates used inside our comb tests and our kinematic-feature control.
